\section{Models and common preliminaries}\label{sec:model}

The four models differ in small but consequential ways.  This section
collects their common building blocks so that the later arguments need not
repeat them.  We first define the models and adversaries, then develop the
offline pair identities and fluid completion curves, and finally record
the probabilistic and finite-size tools used throughout the paper.

\subsection{Models}

There are $n$ jobs, all available at time zero, and one machine.  Operations
are nonpreemptive: once an operation starts, it runs to completion.  At most
one operation can run at a time.  Job $j$ has
a hidden value $p_j\ge0$.  The four models differ in what can be done before
this value is known.

\begin{description}[leftmargin=*,style=nextline]
\item[Obligatory testing $\mathsf{OT}$.]
The values $p_j$ are arbitrary finite nonnegative numbers.  Every job must
first be tested for one unit.  The test reveals $p_j$, after which its
processing operation of length $p_j$ may be performed immediately or
delayed.

\item[Blind execution $\mathsf{BE}$.]
The values $p_j$ are arbitrary finite nonnegative numbers.  A job may be run
blindly for its actual duration $p_j$;
the value becomes known only when that execution finishes.  Alternatively,
the job may be tested for one unit, revealing $p_j$, and its processing
operation may then be delayed.  Testing provides information but does not
shorten the job.

\item[Blind optimization $\mathsf{BO}(u)$.]
Here $p_j\in[0,u]$.  A job may be run raw for the known time $u$.
Alternatively, one unit of optimization changes its running time from $u$
to $p_j$ without revealing $p_j$; the optimized job is subsequently run
blindly, possibly after a delay.  Its value becomes observable only when that
execution finishes.  Raw execution reveals no hidden value.

\item[Revealing optimization $\mathsf{RO}(u)$.]
Again $p_j\in[0,u]$, and raw execution takes the known time $u$.  The other
option is a unit test that reveals $p_j$, followed, possibly after a delay,
by processing for $p_j$.  Thus the preliminary action provides both the
speedup and information about its size.
\end{description}

The notation $\mathsf{BE}$ has no model parameter: blind execution permits
arbitrary finite nonnegative processing times.  For its bounded
instance-optimal theorem we restrict the inputs to $[0,L]$ and make $L$
public.  In the last two models, by contrast, $u$ is itself an available
running time and hence part of every scheduling decision.  We call it the
\emph{raw-execution cap}, or simply the \emph{cap}.  Since the offline
effective length in these models is $\min\{u,1+p_j\}$, we call it the
\emph{offline cap $u$} when discussing the clairvoyant benchmark.  The
number $n$
and any theorem-specific bound $L$ or model parameter $u$ are known to the
algorithm.

The completion time $C_j$ is the time at which job $j$ finishes.  If a test
reveals $p_j=0$, the job completes at the end of that test.  A blind zero job
in $\mathsf{BE}$ completes as soon as it is selected, whereas an optimized
zero job in $\mathsf{BO}(u)$ first pays its unit optimization cost.  The
objective is
\[
                              \sum_{j=1}^n C_j.
\]

Obligatory testing can be viewed as the limit of revealing optimization in
which raw execution becomes unavailable.  We use this only as intuition.
The obligatory theorems and their algorithms are proved directly; deriving
them from revealing optimization would require arguing about both $u$ and $n$ going to infinity, which we avoid.

We write $\ALG_{\mathcal A}(I)$ for the objective value of an online
algorithm $\mathcal A$ on an instance $I$, and $\OPT(I)$ for the value of an
offline optimum that knows all processing times.  We omit algorithm
subscripts or instance arguments when they are clear.
When a sequence of costs has the form $cn^2+o(n^2)$, we call $c$ its
\emph{leading coefficient}; terms contributed by individual jobs rather
than pairs are typically only $O(n)$.

\subsection{Input, information and adversaries}
\label{par:input-convention}

For a deterministic algorithm, the adversary may choose the processing-time
vector and its assignment to labels with full knowledge of the algorithm.
The randomized results use an \emph{oblivious} adversary: the processing-time
vector is fixed before the algorithm draws its private random seed.

An input is a vector $p=(p_1,\ldots,p_n)$ of job lengths chosen by the
adversary.  For a randomized algorithm, this vector is then relabeled by a
uniform permutation.  We abbreviate the expectation over this shuffle and
all further private randomness by $\EE\ALG_{\mathcal A}(p)$.  Thus randomized
performance depends only on the entries of $p$, not on their displayed
order.  An \emph{announced} algorithm is additionally given the multiset
$\{p_1,\ldots,p_n\}$, but not its shuffled assignment to labels.
Every algorithm is \emph{nonanticipating}: each action may depend on the
public parameters, its private randomness, and everything revealed so far,
but not on a hidden value that the model has not yet exposed.
The \emph{transcript} of a run is the resulting ordered sequence of actions
and observations.

For some lower bounds it is convenient to describe a finite distribution
over inputs before selecting one fixed input.  If every deterministic
algorithm has average payoff at least $L$ under this distribution, then the
same is true after conditioning on the private seed of a randomized
algorithm.  Averaging over that seed shows that at least one fixed input in
the support gives the randomized algorithm expected payoff at least $L$.
This is the finite one-sided form of Yao's principle~\cite{Yao1977}.

We write
\begin{equation}
 D[p]\defeq\frac1n\sum_{i=1}^n\delta_{p_i}
 \label{eq:empirical-distributions}
\end{equation}
for the empirical distribution of the input, where $\delta_x$ denotes one
unit of probability concentrated at $x$.  It is unchanged by the
shuffle.

Our instance-optimality results compare against all algorithms receiving
this shuffled input.  This is ordinary instance optimality when the input is
viewed as a multiset, and order-oblivious instance optimality when it is
viewed as an ordered vector; see
\cite{AfshaniBarbayChan2017,NarayananRozhonTetekThorup2024}.

\subsection{Offline optima and pair accounting}

The offline optimum knows all values $p_j$.  It therefore chooses the shorter
available way to complete each job and orders the resulting job blocks
shortest first.  Associate with job $j$ the model-dependent
\emph{effective length}
\begin{equation}
 \lambda_j\defeq
 \begin{cases}
  1+p_j, & \mathsf{OT},\\
  p_j, & \mathsf{BE},\\
  \min\{u,1+p_j\}, & \mathsf{BO}(u)\text{ or }\mathsf{RO}(u).
 \end{cases}
 \label{eq:model-effective-length}
\end{equation}

\begin{lemma}[Effective lengths and the offline pair identity]
\label{lem:offline}\label{lem:optional-offline}
Fix any one of the four models and let
$\lambda_{(1)}\le\cdots\le\lambda_{(n)}$ be its effective lengths from
\cref{eq:model-effective-length}.  A clairvoyant (fully informed) optimum
completes the jobs
in this order and has value
\begin{align}
 \OPT
 &=\sum_{k=1}^n\sum_{i=1}^k\lambda_{(i)},
 \label{eq:unified-opt-prefix}\\
 &=\sum_i\lambda_i+\sum_{i<j}\min\{\lambda_i,\lambda_j\},
 \label{eq:optional-opt-pair}\\
 &=\frac12\left(
      \sum_i\lambda_i+\sum_{i,j}\min\{\lambda_i,\lambda_j\}
    \right).
 \label{eq:unified-opt-symmetric}
\end{align}
In $\mathsf{OT}$ it tests every job.  In $\mathsf{BE}$ it runs every job
blindly.  In $\mathsf{BO}(u)$ and $\mathsf{RO}(u)$ it uses raw mode when
$u\le1+p_j$ and the unit preliminary action followed by processing otherwise.
\end{lemma}

\begin{proof}
Let $C_{[1]}\le\cdots\le C_{[n]}$ be the ordered completion times of any
feasible schedule.  Completing job $j$ consumes at least $\lambda_j$ units of work
belonging to that job.  Hence, by time $C_{[k]}$, the machine has performed
the required work of some $k$ jobs, and
\[
 C_{[k]}\ge\sum_{i=1}^k\lambda_{(i)}.
\]
Choose the mode specified in the lemma, keep each resulting block contiguous,
and execute the blocks in nondecreasing order of $\lambda_j$.  This schedule is
feasible in the relevant model and attains equality for every $k$.  Summing
over $k$ and then collecting the contribution of each unordered pair gives
the first two formulas.  Separating the terms with $i=j$ from the two
orientations of every pair with $i\ne j$ gives the symmetric form.
\end{proof}

For later pair calculations, if $X\subseteq[n]$, we write
\begin{equation}
 \SPT(X)\defeq
 \sum_{i\in X}p_i
 +\sum_{\substack{i,j\in X\\i<j}}\min\{p_i,p_j\}
 \label{eq:spt-block-cost}
\end{equation}
be the exact cost of processing that block in shortest-processing-time
(SPT) order.  Thus \(\SPT(X)\) is the processing-time part of the offline
pair formula when the jobs in $X$ are executed without preliminary actions.

\subsection{Competitive-ratio conventions}

Every asymptotic statement in this paper allows a uniform additive
$o(n^2)$ term.  Formally, an algorithm family $\mathcal A$ has asymptotic
competitive ratio at most $R$ if there is a deterministic sequence
$\eps(n)\to0$ such that every $n$-job instance satisfies
\begin{equation}
 \EE\ALG_{\mathcal A}(I)
 \le R\cdot\OPT(I)+\eps(n)n^2.
 \label{eq:cr-infty}
\end{equation}
For a deterministic algorithm the expectation is omitted.  The optimal ratio
in a model is the infimum of such $R$ over all admissible algorithm families.
We also call this the \emph{size-asymptotic} ratio to emphasize that the
model parameters are fixed while the number of jobs $n$ tends to infinity.

In obligatory testing every effective length is at least one, and in either
optimization model it is at least $\min\{u,1\}$.  Hence
\cref{eq:unified-opt-prefix} gives $\OPT=\Omega(n^2)$ in these models, and the
definition above is equivalent to the usual worst-case ratio after taking
$n\to\infty$.  This equivalence fails for blind execution, where $\OPT$ can
be $o(n^2)$ in the special case where most of the jobs have length zero.  Accordingly,
\cref{thm:blind-instance} makes the absolute statement
$\EE\ALG=n^2\PhiBE(D[p])+o_L(n^2)$, where the benchmark $\PhiBE$ is defined
in \cref{sec:blind-benchmark}.  It does not claim a bounded ratio to the
clairvoyant optimum when most jobs have length zero.

\subsection{A general fluid schedule and its analysis}
\label{sec:common-fluid-envelopes}

The instance-optimality proofs later in the paper repeatedly analyze the same
schedule.  It is useful to describe that schedule in full before deriving
its cost.  Let $D$ be the distribution of processing times, let $q$ be the
fraction of jobs that we intend to test, and suppose that every untested job
can instead be completed at a common fluid cost $v$.  In revealing
optimization $v=u$, the raw execution time.  In blind execution $v$ is the
mean processing time.

At the fluid scale, the $n$ jobs are replaced by one unit of divisible job
mass: a group containing $30\%$ of the jobs has mass $0.3$.  Work is divided
by $n$ as well, so processing this group for time $p$ uses $0.3p$ units of
work.

\paragraph{General fluid schedule.}
Given $D,q$, and $v$, perform the following steps.

\begin{enumerate}[label=\textnormal{(\roman*)},leftmargin=*]
 \item Reserve a $q$-fraction of every processing-time group for testing and
       assign the remaining $1-q$ fraction to the other completion mode.  In
       the finite online algorithm this split is produced by a private random
       order; it does not use the hidden processing times.
 \item Whenever a tested job reveals a sufficiently short processing time,
       process it together with its test.  Defer the other positive tested
       jobs.  The threshold for ``sufficiently short'' will be denoted by
       $\tau_D$.
 \item Regard the tests together with their short outcomes as one block.
       For every $p$, the deferred jobs of processing time $p$ form a block
       with work $p$ per completion.  The jobs assigned to the other mode
       form one more block with work $v$ per completion.
 \item Execute these blocks in increasing order of work per completed job.
       Within the testing block, each selected short outcome is processed
       directly after its test.
\end{enumerate}

The order in the last step can be stated concretely.  Every deferred job has
processing time at least $\tau_D$.  If $v<\tau_D$, the schedule first uses
the other mode, then runs the testing block, and finally processes all
deferred jobs in SPT order.  If $v\ge\tau_D$, it runs the testing block first,
then the deferred jobs with $p<v$, then the other mode, and finally the
deferred jobs with $p\ge v$.  Ties may be resolved arbitrarily except that
the testing block must precede a deferred block of the same length.

There are two things to prove.  First, we must choose $\tau_D$ so that the
testing block completes as many jobs per unit of work as possible.  Second,
we must prove that, after ordering the blocks shortest first, no more
complicated interleaving can complete jobs earlier.  Once this is known, the
leading total completion time is obtained by integrating the mass of jobs
that this schedule has not yet completed.

\paragraph{Choosing the testing block.}
Let $D$ have finite mean.  Imagine testing one unit of job mass.  If
$g(p)\in[0,1]$ is the fraction of jobs with revealed processing time $p$ that
we process immediately, the testing block completes mass $a_D(g)$ and uses
work $w_D(g)$, where
\begin{equation}
 a_D(g)\defeq D(\{0\})+\int_{(0,\infty)}g(p)\,dD(p),
 \qquad
 w_D(g)\defeq 1+\int_{(0,\infty)}p g(p)\,dD(p).
 \label{eq:common-module-aw}
\end{equation}
The term $1$ in $w_D(g)$ pays for the tests, and zero jobs are included in
$a_D(g)$ because they finish at their tests.  Thus $a_D(g)/w_D(g)$ is the
number of completions per unit of work.  We call this combination of tests
and immediate processing a \emph{test module}.

For example, suppose half the mass has length zero and half has length one.
Tests alone complete mass $1/2$ using work $1$.  Processing the length-one
outcomes as well completes the full mass using work $3/2$, improving the
completion density from $1/2$ to $2/3$.  The same choice would be harmful if
the positive outcomes were very long.  In general, an outcome of length $p$
should be included exactly when $p$ is below the current work per completion.
At the optimum, this comparison gives a self-consistent threshold $\tau_D$.

\begin{lemma}[Maximum-density test module]
\label{lem:maximum-density-module}
There is a unique $\tau_D\ge1$ satisfying
\begin{equation}
 \int_{[0,\infty)}(\tau_D-p)^+\,dD(p)=1.
 \label{eq:common-threshold-equation}
\end{equation}
Choose $g_*$ to select every positive $p<\tau_D$, reject every $p>\tau_D$,
and select any fraction of the jobs with $p=\tau_D$.  Write
$a_*=a_D(g_*)$ and $w_*=w_D(g_*)$.  Then
\begin{equation}
 \frac{a_*}{w_*}=\max_g\frac{a_D(g)}{w_D(g)}=\frac1{\tau_D},
 \qquad
 \tau_D a_D(g)\le w_D(g)\quad\text{for every }g.
 \label{eq:common-density-dual}
\end{equation}
\end{lemma}

\begin{proof}
The function $h(t)=\int(t-p)^+\,dD(p)$ is continuous and nondecreasing,
satisfies $h(0)=0$, and tends to infinity because
$h(t)\ge t-\int p\,dD(p)$.  Once positive it
is strictly increasing, so $h(t)=1$ has a unique solution; $h(t)\le t$ gives
$\tau_D\ge1$.

It remains to show that this solution describes the best module.  For an
arbitrary selection rule $g$,
\[
 \begin{aligned}
  \tau_D a_D(g)-\int_{(0,\infty)}p g(p)\,dD(p)
   &=\tau_D D(\{0\})
     +\int_{(0,\infty)}(\tau_D-p)g(p)\,dD(p)\\
   &\le\int_{[0,\infty)}(\tau_D-p)^+\,dD(p)=1.
 \end{aligned}
\]
This is the second inequality in \cref{eq:common-density-dual}.  The threshold
rule $g_*$ attains equality because it keeps exactly the terms with
$\tau_D-p>0$; mass at $p=\tau_D$ does not affect the equality.  Hence
$\tau_D a_*=w_*$ and $g_*$ has maximum density $1/\tau_D$.
\end{proof}

The lemma gives both pieces of information needed later.  The best module
completes mass $a_*$ at work $\tau_Da_*$, and no other combination of tests
and immediate outcomes has smaller work per completion.  The positive jobs
not completed by this module form the residual measure
\begin{equation}
 d\nu(p)=(1-g_*(p))\one_{p>0}\,dD(p),
 \qquad \nu(\{0\})=0.
 \label{eq:common-module-data}
\end{equation}
The measure $\nu$ is not normalized: its total mass $1-a_*$ is exactly the
fraction left for later.  Its support lies in $[\tau_D,\infty)$ because the
threshold rule accepts every shorter outcome.  This ordering fact will make
it possible to execute the test module before all residual work.

For an empirical input, this threshold rule has an elementary finite
interpretation: sort the jobs by processing time and take a prefix.

\begin{corollary}[Empirical maximum-density prefix]
\label{cor:empirical-density-prefix}
For an empirical distribution $D[p]$ from
\cref{eq:empirical-distributions}, the maximum of
\begin{equation}
 \frac{|E|/n}{1+n^{-1}\sum_{i\in E}p_i}
 \label{eq:empirical-prefix-density}
\end{equation}
over sets $E\subseteq[n]$ containing all zero jobs and satisfying
$p_i\le p_j$ whenever $i\in E$ and $j\notin E$ is
$1/\tau_{D[p]}$.  It is attained by taking every occurrence below
$\tau_{D[p]}$, no occurrence above it, and any subset of occurrences equal to
$\tau_{D[p]}$.  Thus every selected value is at most $\tau_{D[p]}$ and every
unselected value is at least $\tau_{D[p]}$.
\end{corollary}

\begin{proof}
Apply \cref{lem:maximum-density-module} to $D[p]$.  At the boundary
$p_i=\tau_{D[p]}$, adding or removing an occurrence preserves the ratio
$1/\tau_{D[p]}$, so the fractional boundary rule may be replaced by any
subset of the tied occurrences.
\end{proof}

\paragraph{Ordering all blocks.}
We now justify step~(iv) of the general schedule.  Besides the test module,
we have the processing of the deferred tested jobs and the other completion
mode.  At the fluid scale, we describe each such block by the mass $c_k$ of
jobs that it can complete and the work $v_k$ required for each of them.  We
write this block as $c_k\delta_{v_k}$.  For example,
$0.3\delta_2$ represents mass $0.3$ that needs two units of work per job.

For a collection
$\mathcal M=\sum_k c_k\delta_{v_k}$ of total mass
$\sum_kc_k=1$, define
\begin{equation}
 K_{\mathcal M}(x)\defeq
 \max\left\{\sum_k z_k:
       0\le z_k\le c_k,\ \sum_kv_kz_k\le x\right\}.
 \label{eq:divisible-spt-envelope}
\end{equation}
This is the largest mass that can be completed using work $x$.  The variables
$z_k$ allow us to use only part of a block.  Thus
$K_{\mathcal M}$ is the \emph{completion curve} of the blocks.

If a finite schedule has completed $N(nx)$ jobs after $nx$ units of actual
work, then its normalized cost is
\[
 \frac{\ALG}{n^2}=\int_0^\infty
       \left(1-\frac{N(nx)}n\right)dx.
\]
The fluid analogue is therefore the area under the unfinished fraction.  For
any completion curve $C$, write
\begin{equation}
 \Area(C)\defeq\int_0^\infty(1-C(x))\,dx
 \label{eq:completion-curve-area}
\end{equation}
for this quantity.  We also write
\begin{equation}
 \SPT(\xi)\defeq
 \frac12\iint\min\{v,w\}\,d\xi(v)d\xi(w).
 \label{eq:divisible-spt-cost}
\end{equation}
The next lemma is simply the shortest-processing-time rule for these
divisible blocks.

\begin{lemma}[Shortest-first order for divisible blocks]
\label{lem:divisible-spt-area}
The curve $K_{\mathcal M}$ processes the items in nondecreasing order of
$v_k$, and its area is the divisible SPT cost:
\begin{equation}
 \Area(K_{\mathcal M})=\SPT(\mathcal M).
 \label{eq:divisible-spt-area}
\end{equation}
\end{lemma}

\begin{proof}
Consider two consecutive blocks with masses $c,c'$ and work per completion
$v\le v'$.  Putting the $v$-block first changes the area, relative to the
opposite order, by $cc'(v-v')\le0$.  Repeating this neighboring swap orders
all blocks by nondecreasing work per completion.  In a
block of completion mass $c_k$ and work per completion $v_k$, unfinished
mass decreases linearly.  After sorting the blocks by $v_k$, its area
contribution is
\[
 v_kc_k\left(1-\sum_{h<k}c_h-\frac{c_k}{2}\right).
\]
Summing these terms and collecting each unordered pair gives
\cref{eq:divisible-spt-area}.
\end{proof}

\paragraph{Why the general schedule is optimal.}
Fix the tested fraction $q$.  We prove a stronger statement than merely
computing the cost of our schedule: after any amount of work, no schedule
that tests mass $q$ can have completed more jobs.

Assume for the moment that $D$ has finite support and put
$d_p=D(\{p\})$, the mass of jobs with processing time $p$.  Consider an
arbitrary schedule after it has performed $x$ units of work.  Let
$t$ be the mass already tested, $b$ the mass completed by the other mode,
and $c_p$ the mass of tested jobs with processing time $p>0$ that has also
been processed.  Thus this schedule has completed mass
\[
 d_0t+b+\sum_{p>0}c_p,
\]
whereas the work used is $t+vb+\sum_{p>0}pc_p$.  We cannot test more than
$q$ mass, use the other mode on more than $1-q$ mass, or process more than
$d_pt$ mass of jobs with processing time $p$.  Consequently, every such
schedule has completed at most the following quantity:
\begin{equation}
 \mathcal C^v_{D,q}(x)\defeq
 \max\left\{
 d_0t+b+\sum_{p>0}c_p:
 \begin{array}{l}
 0\le t\le q,\quad 0\le b\le1-q,\\
 0\le c_p\le d_pt,\\
 t+vb+\sum_{p>0}pc_p\le x
 \end{array}\right\}.
 \label{eq:common-fluid-envelope}
\end{equation}
The function $\mathcal C^v_{D,q}(x)$ is therefore an upper bound on how much
mass any schedule can complete by work $x$.  We next show that the general
schedule described at the start of the subsection attains this bound.  Its
three blocks are
\[
 \underbrace{qa_*\delta_{\tau_D}}_{\substack{\text{tests and selected}\\
                                              \text{short outcomes}}}
 \;+
 \underbrace{q\nu}_{\text{deferred tested outcomes}}
 \;+
 \underbrace{(1-q)\delta_v}_{\text{other mode}}.
\]
The first block has mass $qa_*$ and work $\tau_D$ per completion.  The
second contains the residual processing times themselves.  The last block
has the capacity and cost of the other mode.  The next lemma says that
ordering these blocks shortest first attains $\mathcal C^v_{D,q}(x)$
simultaneously for every $x$.  In particular, the linear program is not
choosing a different schedule for each value of $x$.

\begin{lemma}[Optimality of the general fluid schedule]
\label{lem:test-module-envelope}
Let $\tau_D,g_*$, and $a_*$ be as in
\cref{lem:maximum-density-module}, and let $\nu$ be the deferred mass from
\cref{eq:common-module-data}.  Then
\begin{equation}
 \mathcal C^v_{D,q}=K_{\mathcal M_{D,q,v}},
 \qquad
 \mathcal M_{D,q,v}\defeq
  qa_*\delta_{\tau_D}+q\nu+(1-q)\delta_v.
 \label{eq:common-envelope-items}
\end{equation}
Consequently,
\begin{equation}
 \Area(\mathcal C^v_{D,q})=\SPT(\mathcal M_{D,q,v}),
 \label{eq:common-envelope-area}
\end{equation}
and the general fluid schedule minimizes total completion time among all
schedules that test mass $q$.
\end{lemma}

\begin{proof}
We first show that no feasible schedule can complete more mass than the three
blocks, and then show that the general schedule realizes those blocks.

\emph{Compressing an arbitrary schedule.}
For each processing time $p>0$, call a fraction $g_*(p)$ of its jobs
\emph{selected} and the rest \emph{deferred}; this is only bookkeeping, and
the arbitrary schedule need not respect this division.  For a feasible point
in \cref{eq:common-fluid-envelope}, let $y_{\rm test}$ be the completed mass
coming from tests, tested zeros, and processed selected outcomes, and let
$x_{\rm test}$ be the work spent on them.
Applying \cref{eq:common-density-dual} to the processed fractions within the
first $t$ units of tests gives
\[
 \tau_D y_{\rm test}\le x_{\rm test},
 \qquad y_{\rm test}\le a_*t.
\]
With $\lambda=y_{\rm test}/a_*$, replace this part by $\lambda$ full test modules.
The replacement preserves completed mass, weakly decreases work, and uses
$0\le\lambda\le t\le q$.  Deferred jobs with processing time $p$ have total
capacity at most $q\nu(\{p\})$, while the other completion mode has capacity
$1-q$.  Thus every feasible point gives a feasible choice in the definition
of $K_{\mathcal M_{D,q,v}}$ from
\cref{eq:divisible-spt-envelope,eq:common-envelope-items}.

\emph{Realizing the three blocks.}
Order the three kinds of blocks by work per completion,
placing the test module before a deferred block in case of a tie.  The
other-mode block can be run whenever this order requires it.  A fraction of
the test-module block is performed by running the corresponding fraction of
tests and processing their selected short outcomes.  Every deferred
processing time is at least $\tau_D$, so the test module precedes every
deferred block.  Hence all deferred jobs have been exposed before the
schedule processes them.  The schedule is feasible and realizes
$K_{\mathcal M_{D,q,v}}(x)$ at every work budget $x$.  Its area is given by
\cref{lem:divisible-spt-area}.
\end{proof}

\paragraph{Rounding processing times.}
The later finite algorithms round processing times to a grid.  Moving every
positive value by at most $h$ shifts any fluid completion curve horizontally
by at most $h$, and hence changes its area by at most $h$.  The only caveat is
zero: a tested zero completes without a processing operation, so rounding
must preserve exactly which jobs have value zero.

\begin{lemma}[Stability under zero-preserving rounding]
\label{lem:zero-preserving-transport}
Suppose there are joint random variables $(P,P')$, with respective laws
$D,D'$, such that
$P=0$ if and only if $P'=0$ and $|P-P'|\le h$ almost surely.  Transporting
the same fluid algorithm through this coupling preserves all completion events
and changes accumulated processing work by at most $h$.  Consequently, each
model-specific fluid value $\Phi$ defined below satisfies
\begin{equation}
 |\Phi(D)-\Phi(D')|\le h.
 \label{eq:zero-preserving-transport}
\end{equation}
\end{lemma}

\begin{proof}
Pair each infinitesimal job occurrence under the stated joint distribution
and make the same decision for the two paired values, exposing them only
when the model exposes the corresponding processing time.  Zero preservation
makes test completions coincide.  At
every prefix, the total mass that has received processing is at most one, so
replacing its processing times changes cumulative work by at most $h$.  One
unfinished-mass curve is therefore contained in a horizontal $h$-shift of
the other, and its area changes by at most $h$.  Taking the infimum over
algorithms gives one inequality; exchanging $D,D'$ gives the other.
\end{proof}

\paragraph{Learning from a pilot sample.}
The instance-optimal algorithms do not know $D$.  They use a small pilot
sample with empirical law $\widehat D$, choose the best fluid parameters for
that estimate, and run the resulting schedule on the remaining jobs, whose
empirical law we denote by $D^{\rm rem}$.  We call a fixed choice of these
few parameters---for example, a threshold and a tested fraction---a
\emph{template} $\pi$, and write $F_D(\pi)$ for its fluid cost on $D$.

The learning argument used later is elementary.  Suppose every template is
$L_0$-Lipschitz in a distance $d$, and let $\widehat\pi$ minimize
$F_{\widehat D}$.  Comparing the same template first on $D^{\rm rem}$ and
$\widehat D$, then comparing their respective minimizers, gives
\[
 F_{D^{\rm rem}}(\widehat\pi)
 \le \Phi(D^{\rm rem})
      +2L_0d(\widehat D,D^{\rm rem}).
\]
Comparing the remaining input with the original one gives
\[
 \Phi(D^{\rm rem})
 \le \Phi(D)+L_0d(D^{\rm rem},D).
\]
Thus pilot learning costs two copies of the estimation error and one copy of
the change caused by removing the pilot.  In each application we combine
these two inequalities directly with the model-specific finite-schedule
error.

\subsection{Probabilistic tools}
\label{sec:probabilistic-tools}

We collect here the concentration estimates used in the randomized sections.
All of them are finite-space statements.  We use without further mention the
union bound, Markov's inequality
$\Prob(Y\ge t)\le \EE Y/t$ for $Y\ge0$, and the tail identity
$\EE Y=\int_0^\infty\Prob(Y\ge t)\,dt$.

\paragraph{Deferred decisions.}
Suppose the coordinates of a fixed vector $p$ are assigned uniformly to the
public job labels by the private shuffle.  Even if an algorithm chooses the next untouched label from
everything it has seen so far, the processing time behind that label is
uniform among the remaining occurrences.  Repeating this observation shows
that the processing times encountered at first touch form a uniform random
permutation of the coordinates of $p$.  Equal values may be regarded as distinct
copies when making this argument.

The action chosen at a first touch is determined before the processing time
behind that label is exposed.  We call any choice made from past observations
but before the next value is seen \emph{predictable}.  Such an action can be
used as a selector in the sampling bounds below.  If the chosen action does not reveal
the processing time to the algorithm, we may still expose it afterwards for the
analysis; this does not change the algorithm's decisions.

We next record the concentration bounds applied to the resulting random
orders.  Chernoff bounds handle the few places where genuinely independent
Bernoulli variables arise.

\begin{lemma}[Chernoff bounds]
\label{lem:chernoff}
Let $X=\sum_{i=1}^m X_i$, where the $X_i$ are independent Bernoulli random
variables, and write $\mu=\EE X$.  For every $\delta>0$,
\begin{equation}
 \Prob\bigl(X\ge(1+\delta)\mu\bigr)
 \le \left(\frac{e^\delta}{(1+\delta)^{1+\delta}}\right)^\mu
 \le \exp\!\left(-\frac{\delta^2\mu}{2+\delta}\right).
 \label{eq:chernoff-upper}
\end{equation}
For $0\le\delta\le1$,
\begin{equation}
 \Prob\bigl(X\le(1-\delta)\mu\bigr)
 \le \exp\!\left(-\frac{\delta^2\mu}{2}\right).
 \label{eq:chernoff-lower}
\end{equation}
\end{lemma}

These are the standard exponential-moment bounds of
Chernoff~\cite{Chernoff1952}.

Most of our samples are drawn from a fixed input without replacement.  The
next lemma gives both a tail bound and the finite-population variance used
for such samples.

\begin{lemma}[Sampling without replacement]
\label{lem:sampling-without-replacement}
Let $x_1,\ldots,x_N\in[0,1]$, let $X_1,\ldots,X_N$ be a uniformly random
permutation of these numbers, and put
$\bar x=N^{-1}\sum_i x_i$.  For $1\le k\le N$ and $t\ge0$,
\begin{equation}
 \Prob\left(\left|\sum_{i=1}^kX_i-k\bar x\right|\ge t\right)
 \le 2\exp\!\left(-\frac{2t^2}{k}\right).
 \label{eq:without-replacement-hoeffding}
\end{equation}
Moreover, if $\widehat x_k=k^{-1}\sum_{i=1}^kX_i$ and
$\sigma^2=N^{-1}\sum_i(x_i-\bar x)^2$, then, for $N>1$,
\begin{equation}
 \operatorname{Var}(\widehat x_k)
 =\frac{N-k}{k(N-1)}\sigma^2
 \le\frac{1}{4k}.
 \label{eq:without-replacement-variance}
\end{equation}
\end{lemma}

\begin{proof}
The estimate \cref{eq:without-replacement-hoeffding} is Hoeffding's
sampling-without-replacement inequality~\cite[Section~6]{Hoeffding1963}.
The variance identity is the elementary finite-population correction.
\end{proof}

Applying the scalar variance bound to class indicators gives a concise
estimate for an empirical histogram.

\begin{lemma}[Histogram error without replacement]
\label{lem:without-replacement-histogram}
Partition a population of $N$ occurrences into $d$ classes.  Let $p$ be the
vector of class frequencies in the population and let $\widehat p$ be the
corresponding frequency vector among the first $k$ occurrences in a uniformly
random permutation.  Then
\begin{equation}
 \EE\|\widehat p-p\|_1\le\sqrt{\frac d k}.
 \label{eq:without-replacement-histogram}
\end{equation}
\end{lemma}

\begin{proof}
For a class of mass $p_j$, \cref{eq:without-replacement-variance} and
Cauchy--Schwarz give
$\EE|\widehat p_j-p_j|\le\sqrt{p_j/k}$.  Summing over the classes and using
$\sum_j\sqrt{p_j}\le\sqrt d$ proves the claim.
\end{proof}

An adaptive algorithm may stop at a data-dependent time.  We therefore also
need one event that controls every class at every possible prefix.

\begin{lemma}[Uniform prefix histogram bounds]
\label{lem:uniform-prefix-histogram}
In the same population, let $H_{k,j}$ count class $j$ among the first $k$
entries of a uniformly random permutation.  Then
\begin{align}
 \Prob\left(
   \max_{0\le k\le N}\max_{1\le j\le d}|H_{k,j}-kp_j|\ge t
 \right)
 &\le 2d(N+1)\exp\!\left(-\frac{2t^2}{N}\right),
 \label{eq:uniform-prefix-hoeffding}\\
 \EE\max_{0\le k\le N}\max_{1\le j\le d}|H_{k,j}-kp_j|
 &\le C\sqrt{N\ln(2d(N+1))}
 \label{eq:uniform-prefix-mean}
\end{align}
for a universal constant $C$.
\end{lemma}

\begin{proof}
Apply \cref{eq:without-replacement-hoeffding} to every binary class indicator
and union-bound over $k$ and $j$ to obtain
\cref{eq:uniform-prefix-hoeffding}.  Integrating that tail bound proves
\cref{eq:uniform-prefix-mean}.
\end{proof}

Predictable sampling introduces dependent increments rather than a fixed
prefix.  The following maximal inequalities provide the needed uniform
control.

\begin{lemma}[Martingale maximal inequalities]
\label{lem:martingale-maximal}
Let $(M_k,\mathcal F_k)_{k=0}^N$ be a real martingale with $M_0=0$; that is,
$\mathcal F_k$ records the information available by time $k$ and
$\EE[M_k\mid\mathcal F_{k-1}]=M_{k-1}$.
If $\EE|M_N|^2<\infty$, then Doob's inequality gives
\begin{equation}
 \EE\max_{0\le k\le N}|M_k|^2\le4\EE|M_N|^2.
 \label{eq:doob-L2}
\end{equation}
If instead $|M_k-M_{k-1}|\le b_k$ almost surely for deterministic $b_k$,
then maximal Azuma--Hoeffding gives, for every $t\ge0$,
\begin{equation}
 \Prob\left(\max_{0\le k\le N}|M_k|\ge t\right)
 \le2\exp\!\left(-\frac{t^2}{2\sum_{k=1}^N b_k^2}\right).
 \label{eq:maximal-azuma}
\end{equation}
\end{lemma}

We use the classical forms due to Doob~\cite{Doob1953} and
Azuma~\cite{Azuma1967}; the maximal version of the latter follows by stopping
the usual exponential bound when the displayed threshold is first crossed.

The next lemma models an adaptive algorithm reading a randomly permuted fixed
input.  The vector $X_k$ records the statistics of the $k$th touched job,
while the predictable weight $A_k$ records whether the algorithm selects that
job before seeing its value.  The lemma says that such adaptive selection
cannot systematically bias the sample: every selected prefix remains close
to the corresponding population average, simultaneously over all possible
stopping times.  We stop before the final $\delta N$ jobs because averages in
the remaining population become unstable near the very end.

\begin{lemma}[Predictable sampling from a random permutation]
\label{lem:predictable-permutation-sampling}
Let $X_1,\ldots,X_N$ be a uniform random permutation of fixed vectors
$x_1,\ldots,x_N\in[0,B]^d$, and let
$\bar x=N^{-1}\sum_i x_i$.  Suppose $A_k\in[0,1]$ is chosen using only the
previous draws, before $X_k$ is exposed.  Then for
$\delta,\eta\in(0,1)$, with probability at least
$1-\eta$, simultaneously for every coordinate $\ell$ and every
$m\le(1-\delta)N$,
\begin{equation}
 \left|\sum_{k\le m}A_k(X_{k,\ell}-\bar x_\ell)\right|
 \le(1+\delta^{-1})E_{N,d,B,\eta},
 \qquad
 E_{N,d,B,\eta}\defeq
 B\sqrt{2N\ln\!\left(\frac{4d(N+1)}{\eta}\right)}.
 \label{eq:predictable-permutation-sampling}
\end{equation}
The statement remains valid after conditioning on independent private
randomness used to choose the predictable sequence.
\end{lemma}

\begin{proof}
For each coordinate, apply
\cref{eq:without-replacement-hoeffding} after scaling by $B$, and union-bound
over coordinates and prefix lengths.  Except with probability at most
$\eta/2$, every centered prefix sum has absolute value at most
$E=E_{N,d,B,\eta}$.  Before draw $k\le(1-\delta)N$, the mean
$\widetilde x_{k-1,\ell}$ of
the remaining population therefore satisfies
\[
 |\widetilde x_{k-1,\ell}-\bar x_\ell|\le \frac{E}{\delta N}.
\]
For each $\ell$,
\[
 M_{m,\ell}=\sum_{k\le m}A_k(X_{k,\ell}-\widetilde x_{k-1,\ell})
\]
is a martingale with increments bounded by $B$.  The maximal
Azuma--Hoeffding inequality \cref{eq:maximal-azuma}, followed by a union
bound over $\ell$, puts all $|M_{m,\ell}|$ below $E$ except with probability at
most $\eta/2$.  The accumulated difference between the remaining-population
mean and the original mean is at most $E/\delta$, which proves
\cref{eq:predictable-permutation-sampling}.  Conditioning and then averaging
proves the final assertion.
\end{proof}
